\section{A Subgraph That Owns Nothing}
\label{sec:ch09-a-subgraph-that-owns-nothing}

\index{Ratings service!the second shape of seam|(}%
The Speakers service is the easy seam. A type moved, whole, and the service it
moved to owns it the way the Sessions service owns a session. The Ratings
service is the other shape, and everything awkward in this chapter is in it: it
adds fields to \code{Session}, and it will never own a session.

\index{ownership!a service can own a field without owning the type}%
That is a real arrangement and not a contrivance. Somebody has to store
the scores attendees leave, nobody wants that table inside the service that
owns the program, and a client asking for a session's average score should
not have to know that two systems were involved. Federation's answer is that a
type is the union of what every subgraph says about it, and this is the first
time this book has needed the second half of that sentence.

\subsection{The Rows It Does Own}

\index{Ratings service!the project file}%
The project file is the Speakers one with a different name, and the launch
profile pins the third port:

\begin{minted}[fontsize=\footnotesize]{xml}
<Project Sdk="Microsoft.NET.Sdk.Web">

  <PropertyGroup>
    <TargetFramework>net10.0</TargetFramework>
    <Nullable>enable</Nullable>
    <ImplicitUsings>enable</ImplicitUsings>
    <GenerateDocumentationFile>true</GenerateDocumentationFile>
    <NoWarn>$(NoWarn);CS1591</NoWarn>
  </PropertyGroup>

  <ItemGroup>
    <PackageReference Include="HotChocolate.ApolloFederation" Version="16.6.1" />
    <PackageReference Include="HotChocolate.AspNetCore" Version="16.6.1" />
    <PackageReference Include="HotChocolate.AspNetCore.CommandLine" Version="16.6.1" />
    <PackageReference Include="HotChocolate.Data.EntityFramework" Version="16.6.1" />
    <PackageReference Include="HotChocolate.Types.Analyzers" Version="16.6.1" />
    <PackageReference Include="Microsoft.EntityFrameworkCore.Sqlite" Version="10.0.11" />
  </ItemGroup>

</Project>
\end{minted}

\begin{minted}{json}
{
  "$schema": "https://json.schemastore.org/launchsettings.json",
  "profiles": {
    "http": {
      "commandName": "Project",
      "dotnetRunMessages": true,
      "launchBrowser": false,
      "applicationUrl": "http://localhost:5003",
      "environmentVariables": {
        "ASPNETCORE_ENVIRONMENT": "Development"
      }
    }
  }
}
\end{minted}

\begin{minted}{csharp}
namespace Ratings;

/// <summary>
/// One score somebody left for one session.
/// </summary>
/// <param name="Id">Stable identifier for this rating.</param>
/// <param name="SessionId">
/// Which session it is for. A number this service stores and cannot
/// check: the sessions are another service's rows.
/// </param>
/// <param name="Score">One to five.</param>
public sealed record Rating(int Id, int SessionId, int Score);
\end{minted}

\index{foreign key!across a seam, no database enforces it}%
\code{SessionId} is a foreign key that no database will ever enforce, and the
comment says so because a reader typing this out deserves to know it is not an
oversight. Referential integrity stops at the process boundary. If a session is
deleted in the other service, these rows point at nothing, and no database
will tell anybody. Somebody has to.

\begin{minted}{csharp}
namespace Ratings;

/// <summary>
/// The scores this service was left with. Session 4 has none, which is
/// the case worth seeding deliberately: an average over no rows is not
/// zero, and the schema has to say which.
/// </summary>
public static class RatingData
{
    public static readonly IReadOnlyList<Rating> Ratings =
    [
        new Rating(1, 1, 5),
        new Rating(2, 1, 4),
        new Rating(3, 1, 3),
        new Rating(4, 2, 4),
        new Rating(5, 2, 5),
        new Rating(6, 3, 3)
    ];
}
\end{minted}

\begin{minted}{csharp}
using Microsoft.EntityFrameworkCore;

namespace Ratings;

/// <summary>
/// This service's own database. It holds ratings and nothing else: no
/// session table, no speaker table, and no foreign key that any
/// database could enforce, because the rows SessionId points at are
/// three miles away in another process.
/// </summary>
public sealed class RatingContext(
    DbContextOptions<RatingContext> options) : DbContext(options)
{
    public DbSet<Rating> Ratings => Set<Rating>();

    protected override void OnModelCreating(ModelBuilder model)
        => model.Entity<Rating>().HasData(RatingData.Ratings);
}
\end{minted}

\index{DataLoader!grouped, because one session has many ratings}%
The loader groups where the other two key one to one, which is the only structural
difference from the two the book already has:

\begin{minted}{csharp}
using Microsoft.EntityFrameworkCore;

namespace Ratings;

public static class RatingDataLoader
{
    /// <summary>
    /// Every rating left for each session in one batch, grouped by
    /// session. A session nobody rated is absent from the dictionary
    /// rather than present with an empty list, which is the difference
    /// between an average of null and an average of zero.
    /// </summary>
    [DataLoader]
    public static async Task<Dictionary<int, Rating[]>>
        GetRatingsBySessionIdAsync(
            IReadOnlyList<int> sessionIds,
            RatingContext db,
            CancellationToken ct)
    {
        var rows = await db.Ratings
            .Where(r => sessionIds.Contains(r.SessionId))
            .ToListAsync(ct);

        return rows
            .GroupBy(r => r.SessionId)
            .ToDictionary(g => g.Key, g => g.ToArray());
    }
}
\end{minted}

\index{Query@\code{Query}!a service that mostly answers about somebody else}%
This service needs a root field for the same reason the last one did, and it
has almost nothing to put there. Everything interesting it knows is about a
type it does not own, so what is left for \code{Query} is a total:

\begin{minted}{csharp}
using Microsoft.EntityFrameworkCore;

namespace Ratings;

[QueryType]
public static class Query
{
    /// <summary>
    /// How many scores this service holds, across every session.
    /// </summary>
    public static async Task<int> GetRatingCountAsync(
        RatingContext db,
        CancellationToken ct)
        => await db.Ratings.CountAsync(ct);
}
\end{minted}

\index{StatementLog@\code{StatementLog}!the third copy}%
\code{StatementLog} from chapter~\ref{ch:data-without-n-plus-1} is here too,
identical apart from its namespace and now on its third copy in this book. The
counts in section~\ref{sec:ch09-the-first-query-across-two} are measured per
service, and a service with no interceptor of its own would have contributed a
zero that meant nothing.

\subsection{Declaring Somebody Else's Type}

\index{Session@\code{Session}!declared a second time, with three fields}%
This service needs a \code{Session} class, and it is not the one the Sessions
service has. It is three fields, and each of them is there for a different
reason:

\begin{minted}{csharp}
using HotChocolate.Types.Relay;

namespace Ratings;

/// <summary>
/// A session, as far as this service is concerned.
/// </summary>
/// <param name="Id">Stable identifier for this session.</param>
/// <param name="Key">The number inside that identifier.</param>
/// <param name="Title">The line that goes in the program.</param>
// Declared here rather than referenced from the Sessions project. Two
// services that would be separate repositories owned by separate teams
// do not share a model assembly, and this book would be demonstrating
// the opposite of what it says if they did.
//
// Three fields where the Sessions service has one, and the split is
// the whole lesson of the seam. Id is the global node id, as a string,
// because that is what arrives and what has to go back. Key is the
// integer inside it, which is the only part this service's own tables
// know about, and it is hidden from the schema because no client has
// any use for it. Title is never stored here at all.
public sealed record Session(
    [property: ID] string Id,
    [property: GraphQLIgnore] int Key,
    string Title);
\end{minted}

\index{node id!arrives as a string and has to be split}%
A record that carries the same identifier twice looks wrong and is not. The
key that arrives across the seam is \code{U2Vzc2lvbjox}, and two different
things need two different halves of it: the schema needs the string back, and
the \code{Ratings} table is keyed by the integer inside it. Holding both is
cheaper than decoding twice and safer than storing only the one you happened to
need first.

\subsection{Two Lines in Program.cs, and Why}

\index{AddGlobalObjectIdentification@\code{AddGlobalObjectIdentification}!cannot be used here}%
Every service in this book so far has called
\code{AddGlobalObjectIdentification()}, and this one cannot. Add it and the
service does not start:

\begin{minted}[fontsize=\footnotesize,breakanywhere]{text}
HotChocolate.SchemaException: For more details look at the `Errors` property.

1. There is no object type implementing interface `Node`. (HotChocolate.Types.Relay.NodeType)
   at HotChocolate.SchemaBuilder.Setup.CompleteSchema(...)
\end{minted}

\index{Node@\code{Node}!a subgraph with none cannot switch it on}%
The call adds \code{node}, \code{nodes} and the \code{Node} interface, and then
refuses to build a schema in which nothing implements that interface. This
service owns no node. Its \code{Session} is a contribution, not something it
can fetch by identifier out of a table it has, and giving it a
\code{[NodeResolver]} to satisfy the schema builder would be a lie in a
declaration: it would tell every client that Ratings can answer
\code{node(id:)} for a session, which it cannot.

\index{INodeIdSerializer@\code{INodeIdSerializer}!is registered by the call it cannot make}%
Which leaves a problem, because the keys this service is handed are still
global node ids and it still has to read one. The serializer that reads one is
registered by exactly the call it cannot make. \code{AddDefaultNodeIdSerializer()}
registers it and nothing else, and it is the fifth line of the builder chain in
the file below.

\index{AddDefaultNodeIdSerializer@\code{AddDefaultNodeIdSerializer}!the line that replaces it}%
Leave it out and nothing tells you. The reference resolver is handed a null
serializer, the first call on it throws inside a resolver, and the entity route
answers:

\begin{minted}[fontsize=\footnotesize,breakanywhere]{json}
{"errors":[{"message":"Unexpected Execution Error","path":["_entities"]}],"data":{"_entities":[null]}}
\end{minted}

\index{Unexpected Execution Error@\code{Unexpected Execution Error}!what it means}%
The service builds. The schema builds, and is the document you would have
exported anyway. The process starts, and \code{ratingCount} on its own root
answers correctly at HTTP 200. A smoke test that does not cross the seam
reports a healthy service. Only the entity route fails, and there is no stack
trace, no log line and nothing at startup to say why.

\index{Unexpected Execution Error@\code{Unexpected Execution Error}!means a resolver threw}%
\code{Unexpected Execution Error} is what Hot Chocolate says when a resolver
throws and nothing has been configured to tell you more. It is worth learning to
read as one thing: a resolver asked the container for something and was handed
a null. Chapter~\ref{ch:first-subgraph} met the same message from the same
cause, trying to inject a \code{QueryContext} into a reference resolver.

Here is the whole file:

\begin{minted}{csharp}
using ChilliCream.Nitro.App;
using HotChocolate.AspNetCore;
using Microsoft.EntityFrameworkCore;
using Ratings;

var builder = WebApplication.CreateBuilder(args);

// EF Core reports every command through the logging pipeline as well,
// at Information and with a duration attached. StatementLog below is
// this service's instrument, so silence the other one rather than have
// every statement printed twice.
builder.Logging.AddFilter(
    "Microsoft.EntityFrameworkCore.Database.Command",
    LogLevel.Warning);

builder.Services.AddDbContextFactory<RatingContext>(options =>
{
    options.UseSqlite("Data Source=ratings.db");

    if (builder.Environment.IsDevelopment())
    {
        options.AddInterceptors(new StatementLog());
    }
});

// Two lines here differ from the other two services, and both follow
// from the same fact: this service contributes fields to a type it
// does not own and owns no node of its own.
//
// There is no AddGlobalObjectIdentification, because there cannot be.
// That call refuses to build a schema with no node type in it: "There
// is no object type implementing interface `Node`." So this service
// has no node or nodes field, and needs no interceptor to declare
// them shareable.
//
// But the keys it is handed across the seam are still global node ids,
// and the serializer that reads one is normally registered by the call
// it cannot make. AddDefaultNodeIdSerializer registers it on its own.
// Without this line the reference resolver is injected a null and the
// entity route answers Unexpected Execution Error.
builder
    .AddGraphQL()
    .AddRatingsTypes()
    .RegisterDbContextFactory<RatingContext>()
    .AddDefaultNodeIdSerializer()
    .AddApolloFederation()
    .ModifyCostOptions(options => options.ApplyCostDefaults = false);

var app = builder.Build();

using (var scope = app.Services.CreateScope())
{
    var factory = scope.ServiceProvider
        .GetRequiredService<IDbContextFactory<RatingContext>>();
    using var db = factory.CreateDbContext();
    db.Database.EnsureCreated();
}

app.MapGraphQL()
    .WithOptions(options =>
    {
        // Serve the Nitro build that shipped in the package. The
        // default, ServeMode.Latest, downloads the current one
        // from ChilliCream's CDN instead.
        options.Tool.ServeMode = ServeMode.Embedded;
    });

app.RunWithGraphQLCommands(args);
\end{minted}

\index{ShareNodeFields@\code{ShareNodeFields}!not needed here}%
No \code{ShareNodeFields} either, and for a pleasant reason: with no
\code{node} and no \code{nodes} there is nothing to declare shareable, so the
file chapter~\ref{ch:composition} spent a section justifying is not in this
service at all. Two of the three subgraphs need it and the third does not,
which is a better argument for keeping it out of a shared library than
anything I said in section~\ref{sec:ch09-the-service-that-owns-the-speakers}.

\subsection{The Key Nobody Decodes For You}

\index{entity route!does not decode a node id, three ways}%
The reference resolver here does the same job as the two in
chapter~\ref{ch:first-subgraph}, and what it cannot delegate needs stating
precisely. Chapter~\ref{ch:first-subgraph} established that the key
arrives at a reference resolver as the undecoded node id string, where the node
resolver beside it is handed the integer, and that \code{[ID<T>]} on the
parameter does not fix that: it yields zero.

\index{ID@\code{[ID<T>]}!on the property, the same silent zero}%
The same is true one place over. Put \code{[ID<Session>]} on the property
rather than on the parameter, and declare the resolver to take an \code{int},
and it compiles, exports \code{id: ID!}, composes, runs, and hands the resolver
a zero for every key it is ever given. A session with three ratings answers
with a count of zero and an average of null, at HTTP 200, with no error key.
It reads exactly like a session nobody rated.

\index{node id!decoding is the resolver's job}%
So the decode is the resolver's, here as everywhere. Here is the whole type,
and the two directives on it are section~\ref{sec:ch09-a-field-built-from-a-title}'s
subject, not this one's:

\begin{minted}[fontsize=\footnotesize]{csharp}
using HotChocolate.ApolloFederation.Resolvers;
using HotChocolate.ApolloFederation.Types;
using HotChocolate.Types.Relay;

namespace Ratings;

[ObjectType<Session>]
[Key("id")]
[ReferenceResolver(
    EntityResolver = nameof(ResolveSessionReference))]
public static partial class SessionType
{
    /// <summary>
    /// The line that goes in the program.
    /// </summary>
    // Owned by the Sessions service. Declared here only so that
    // feedbackUrl below has something to require, and resolved by
    // nobody: the router never asks this service for it.
    [External]
    public static string GetTitle([Parent] Session session)
        => session.Title;

    /// <summary>
    /// The mean score, or null if nobody has rated this session.
    /// </summary>
    public static async Task<double?> GetAverageScoreAsync(
        [Parent(nameof(Session.Key))] Session session,
        IRatingsBySessionIdDataLoader ratingsBySessionId,
        CancellationToken ct)
    {
        var ratings = await ratingsBySessionId
            .LoadAsync(session.Key, ct);

        return ratings is { Length: > 0 }
            ? ratings.Average(r => (double)r.Score)
            : null;
    }

    /// <summary>
    /// How many scores this session has been left.
    /// </summary>
    public static async Task<int> GetRatingCountAsync(
        [Parent(nameof(Session.Key))] Session session,
        IRatingsBySessionIdDataLoader ratingsBySessionId,
        CancellationToken ct)
        => (await ratingsBySessionId.LoadAsync(session.Key, ct))
            ?.Length ?? 0;

    /// <summary>
    /// Where an attendee goes to leave a score for this session.
    /// </summary>
    // The url is built from the session's title, and the title belongs
    // to another service. @requires is what makes the router send it
    // along with the key, and [External] on GetTitle above is what
    // makes that declaration legal.
    [Requires("title")]
    public static string GetFeedbackUrl([Parent] Session session)
        => "https://feedback.example/2026/" + Slug(session.Title);

    /// <summary>
    /// Answers one representation from the entity route.
    /// </summary>
    // The key arrives as the undecoded node id string, exactly as it
    // does in the Sessions service, and for the same reason: nothing
    // on the entity route decodes one for you. Marking the property
    // [ID<Session>] does not do it either: it hands this method a zero
    // rather than raising anything, so the service answers with a
    // session that has no ratings and never says why.
    //
    // title is null on every request that does not ask for
    // feedbackUrl, because the router only sends a required field when
    // the field requiring it was asked for.
    [GraphQLIgnore]
    public static Session? ResolveSessionReference(
        string id,
        string? title,
        INodeIdSerializer serializer)
    {
        var nodeId = serializer.Parse(id, typeof(int));

        if (nodeId.TypeName != nameof(Session))
        {
            return null;
        }

        return new Session(
            id,
            (int)nodeId.InternalId!,
            title ?? string.Empty);
    }

    private static string Slug(string title)
        => string.Concat(
            title.ToLowerInvariant()
                .Select(c => char.IsLetterOrDigit(c) ? c : '-'));
}
\end{minted}

\index{type name guard!in a third service}%
The guard on \code{nodeId.TypeName} is chapter~\ref{ch:first-subgraph}'s,
unchanged, in a third service. It matters more here than it did there, because
this service is now genuinely reachable by a router that handles two types of
node id, and a \code{Speaker} key arriving under a \code{Session} type name
would decode to a perfectly good integer and select somebody's ratings.

\index{Parent@\code{[Parent]}!names the hidden key, not the public one}%
On the two resolvers that touch the database, the \code{[Parent]} attribute
names \code{Key} rather than \code{Id}, so what they are handed is the integer
and not the node id. That is the second place in this chapter where the
attribute is load-bearing and easy to miss.

\index{averageScore@\code{averageScore}!null is not zero}%
One design decision in \code{GetAverageScoreAsync} is worth defending, because
the alternative is one character shorter. A session nobody has rated returns
null and not zero. Zero is a score somebody could have given, and a schema
that cannot tell \enquote{nobody rated this} from \enquote{everybody hated
this} has thrown away the difference at the type level, where no client can get
it back. \code{ratingCount} beside it is a genuine zero, and the pair is what
lets a client render \enquote{not yet rated} rather than a one-star average.

\index{parameter binding!a resolver reads the representation by name}%
The second parameter of the reference resolver is
section~\ref{sec:ch09-a-field-built-from-a-title}'s subject. What it shows here
is the binding rule: a reference resolver's parameters are matched against the
fields of the representation by name, so \code{title} arrives if the router
sent one and is null if it did not.

\subsection{The Third Entry}

\index{graph.yaml@\code{graph.yaml}!reaches its final shape}%
One more line in the graph file and the example is complete. Here it is again,
with what chapter~\ref{ch:composition} started, what
section~\ref{sec:ch09-the-first-query-across-two} added, and what this section
adds, all three:

\begin{minted}[highlightlines={11,12,13,14}]{yaml}
version: 1
subgraphs:
  - name: sessions
    routing_url: http://localhost:5001/graphql
    schema:
      file: ../src/Sessions/schema.graphql
  - name: speakers
    routing_url: http://localhost:5002/graphql
    schema:
      file: ../src/Speakers/schema.graphql
  - name: ratings
    routing_url: http://localhost:5003/graphql
    schema:
      file: ../src/Ratings/schema.graphql
\end{minted}

\index{subgraphId@\code{subgraphId}!is positional and renumbers}%
Compose that, and the identifiers chapter~\ref{ch:enter-the-router} taught you
to read out of a query plan are assigned in the order the file lists them:
\code{sessions} is 0, \code{speakers} is 1, \code{ratings} is 2. They are
positional rather than derived from the name, so a subgraph inserted in the
middle of this file renumbers every one after it, and a plan you saved
yesterday names different services today. Worth knowing before you sort the
list alphabetically.
\index{Ratings service!the second shape of seam|)}%
